% Typography follows a continuous technical essay, with numbered equations
% and vector figures. Avoid chapter-per-page or presentation-card layouts.
\usepackage{fontspec}
\setmainfont{Times New Roman}
\setsansfont{Arial}
\setmonofont{Menlo}[Scale=0.88]
\setCJKmainfont{Songti SC}[Script=Default,BoldFont={Songti SC Bold}]
\setCJKsansfont{Heiti SC}[Script=Default]
\setCJKmonofont{Heiti SC}[Script=Default]

\usepackage[left=24mm,right=24mm,top=23mm,bottom=24mm,headheight=16pt]{geometry}
\usepackage{amsmath,amssymb,mathtools}
\usepackage{booktabs,tabularx,array}
\usepackage{graphicx,xcolor}
\usepackage{tikz}
\usetikzlibrary{patterns,arrows.meta,calc}
\usepackage{caption}
\usepackage{needspace}
\usepackage{placeins}
\usepackage{fancyhdr}
\usepackage{hyperref}
\usepackage{bookmark}

\definecolor{ink}{HTML}{18252E}
\definecolor{accent}{HTML}{276481}
\definecolor{selected}{HTML}{DDEBF2}
\definecolor{gridline}{HTML}{ABB8C1}
\definecolor{muted}{HTML}{566571}
\definecolor{statefill}{HTML}{F8E5CB}
\definecolor{stateedge}{HTML}{A5682B}
\definecolor{waste}{HTML}{E9ECEF}
\hypersetup{colorlinks=true,linkcolor=accent,citecolor=accent,urlcolor=accent,
  pdftitle={逐 Query Sparse Attention Kernel 的设计取舍},
  pdfauthor={齐呈祥},pdfsubject={逐 Query Sparse Attention Kernel 的设计取舍}}
\urlstyle{same}
\captionsetup{font=small,labelfont=bf,labelsep=quad,skip=7pt}
\ctexset{
  section={format=\large\sffamily\bfseries,beforeskip=2ex plus .5ex,afterskip=.9ex},
  subsection={format=\normalsize\sffamily\bfseries,beforeskip=1.6ex plus .4ex,afterskip=.7ex}
}
\linespread{1.12}
\raggedbottom
\setlength{\parskip}{2pt plus 1pt}
\clubpenalty=5000
\widowpenalty=5000
\displaywidowpenalty=5000
\renewcommand{\refname}{参考}
\setlength{\emergencystretch}{2em}
\setlength{\textfloatsep}{15pt plus 3pt minus 2pt}
\setlength{\intextsep}{13pt plus 3pt minus 2pt}
\renewcommand{\arraystretch}{1.28}
\newcolumntype{Y}{>{\raggedright\arraybackslash}X}
\pagestyle{fancy}
\fancyhf{}
\fancyhead[LE]{\small\thepage\quad\textbullet\quad 齐呈祥}
\fancyhead[RO]{\small 逐 Query Sparse Attention Kernel 的设计取舍\quad\textbullet\quad\thepage}
\renewcommand{\headrulewidth}{0pt}
\renewcommand{\footrulewidth}{0pt}
\fancypagestyle{plain}{\fancyhf{}\fancyfoot[C]{\small\thepage}\renewcommand{\headrulewidth}{0pt}}

% A compact title lets the opening problem appear on the first page.
\makeatletter
\renewcommand{\maketitle}{%
  \noindent{\LARGE\bfseries\@title\par}
  \vspace{12pt}
  \noindent{\large\@author\par}
  \vspace{8pt}
  \noindent{\small\color{muted}\@date\quad
    \href{https://zhuanlan.zhihu.com/p/2090552659241591553}{知乎原文}\par}
  \vspace{14pt}
}
\makeatother
